\documentclass[12pt]{article}

\usepackage[utf8]{inputenc}
\usepackage{amsmath, amssymb}
\usepackage{geometry}
\geometry{margin=1in}
\usepackage{enumitem}
\setlength{\parindent}{0pt}

\begin{document}
\begin{large}
\textbf{
Problem:}
\end{large}

\vspace{1em}

A worker faces the following income maintenance program:

\begin{itemize}
\item Guaranteed income (benefit): $G = 800$ per week
\item Benefit reduction rate (tax-back rate): $t = 0.4$
\item Market wage: $w = 25$ per hour
\item Total available time: 60 hours per week
\end{itemize}

While eligible for benefits, disposable income is given by:
\[
C = G + (1 - t)wh
\quad \text{for } wh \leq \frac{G}{t}
\]

Once earnings exceed the break-even level, benefits are exhausted and:
\[
C = wh
\]

\vspace{1em}

\textbf{Answer the following questions clearly and show all steps.}

\begin{enumerate}

\item[(a)] \textbf{(3 marks)}  
Calculate:
\begin{enumerate}
\item The break-even level of earnings.
\item The number of hours worked at which benefits are fully exhausted.
\end{enumerate}

\item[(b)] \textbf{(2 marks)}  
Compute:
\begin{enumerate}
\item The effective net wage while the worker is receiving benefits.
\item The effective marginal tax rate on earnings.
\end{enumerate}

\item[(c)] \textbf{(3 marks)}  
Draw the worker’s budget constraint in $(h, C)$ space. Clearly label:
\begin{itemize}
\item The guaranteed income level,
\item The break-even point,
\item The slope in each region.
\end{itemize}

\item[(d)] \textbf{(3 marks)}  
Explain how this program affects:
\begin{enumerate}
\item The incentive to enter the labour force (extensive margin).
\item The incentive to increase hours worked while receiving benefits (intensive margin).
\end{enumerate}
Be precise in distinguishing income and substitution effects.

\item[(e)] \textbf{(4 marks)}  
Suppose the benefit reduction rate increases from $t = 0.4$ to $t = 0.7$, holding all else constant.
\begin{enumerate}[label=(\roman*)]
\item Explain how the budget constraint changes.
\item Predict the likely effect on labour supply.
\item Discuss the equity--efficiency trade-off created by increasing $t$.
\end{enumerate}

\end{enumerate}

\vfill
\textbf{Total: 15 marks}

\end{document}
